\newcommand{\verdicts}{\mathbf{V}}
\newcommand{\xtverdicts}{\mathbf{v}}
\newcommand{\xtculprits}{\mathbf{c}}
\newcommand{\xtfaults}{\mathbf{f}}
\newcommand{\equationtenpointone}{\psi \equiv \tup{\goodset, \badset, \wonkyset, \offenders}}
\newcommand{\equationtenpointtwo}{
    \begin{aligned}
        \xtdisputes \equiv (\xtverdicts, \xtculprits, \xtfaults) \\
        \where \xtverdicts \in \seq{\tup{\H, \ffrac{\tau}{\mathsf{E}} - \N_2, \seq{\tup{\{\top, \bot\}, \N_{\mathsf{V}}, \mathbb{E}}}_{\floor{\nicefrac{2}{3}\mathsf{V}} + 1}}}\\
        \also \xtculprits \in \seq{\H, \H_E, \mathbb{E}} \,,\quad
        \xtfaults \in \seq{\H, \{\top,\bot\}, \H_E, \mathbb{E}}
    \end{aligned}
}
\newcommand{\equationtenpointthree}{
    \begin{align}
        \begin{aligned}
            \forall (r, a, \mathbf{j}) \in \xtverdicts, \forall (v, i, s) \in \mathbf{j} : s \in \sig{\mathbf{k}[i]_e}{\mathsf{X}_v \concat r}\\
            \quad\where \mathbf{k} = \begin{cases}
                \kappa \when a = \displaystyle \ffrac{\tau}{\mathsf{E}}\\
                \lambda \otherwise\\
            \end{cases}
        \end{aligned}\\
    \end{align}
}
\newcommand{\equationtenpointfour}{\mathsf{X}_\top \equiv \text{{\small \texttt{\$jam\_valid}}}\,,\ \mathsf{X}_\bot \equiv \text{{\small \texttt{\$jam\_invalid}}}}
\newcommand{\equationtenpointfive}{
    \begin{align}
        \forall (r, k, s) \in \xtculprits : \bigwedge \left\{
        \begin{aligned}
            r \in \badset' \,,\\
            k \in \mathbf{k} \,,\\
            s \in \sig{k}{\mathsf{X}_G \frown r}
        \end{aligned}\right.\\
    \end{align}
}
\newcommand{\equationtenpointsix}{
    \begin{align}
        \forall (r, v, k, s) \in \xtfaults : \bigwedge \left\{ \begin{aligned}
            r \in \badset' \Leftrightarrow r \not\in \goodset' \Leftrightarrow v \,,\\
            k \in \mathbf{k} \,,\\
            s \in \sig{k}{\mathsf{X}_v \concat r}\\
        \end{aligned}\right.\\
        \nonumber\where \mathbf{k} = \{k_e \mid k \in \lambda \cup \kappa\} \setminus \offenders
    \end{align}
}
\newcommand{\equationtenpointseven}{\xtverdicts = \orderuniqby{r}{\tup{r, a, \mathbf{j}} \in \xtverdicts}}
\newcommand{\equationtenpointeight}{\xtculprits = \orderuniqby{k}{\tup{r, k, s} \in \xtculprits} \,,\quad \xtfaults = \orderuniqby{k}{\tup{r, v, k, s} \in \xtfaults}}
\newcommand{\equationtenpointnine}{\{r \mid \tup{r, a, \mathbf{j}} \in \xtverdicts\} \disjoint \goodset \cup \badset \cup \wonkyset}
\newcommand{\equationtenpointten}{\forall (r, a, \mathbf{j}) \in \xtverdicts : \mathbf{j} = \orderuniqby{i}{\tup{v, i, s} \in \mathbf{j}}}
\newcommand{\equationtenpointeleven}{\verdicts \in \seq{\tup{\H, \{0, \floor{\nicefrac{1}{3}\mathsf{V}}, \floor{\nicefrac{2}{3}\mathsf{V}} + 1\}}}}
\newcommand{\equationtenpointtwelve}{\verdicts = \sq{\tup{r, \sum_{\tup{v, i, s} \in \mathbf{j}}\!\!\!\! v}\ \mid\ \tup{r, a, \mathbf{j}} \orderedin \xtverdicts}}
\newcommand{\equationtenpointthirteen}{\forall (r, \floor{\nicefrac{2}{3}\mathsf{V}} + 1) \in \verdicts : \exists (r, \dots) \in \xtfaults}
\newcommand{\equationtenpointfourteen}{\forall (r, 0) \in \verdicts : |\{(r, \dots) \in \xtculprits\}| \ge 2}
\newcommand{\equationtenpointfifteen}{
    \forall c \in \N_\mathsf{C} : \rho^\dagger[c] = \begin{cases}
        \none \!\!\!\!\when \{ (\mathcal{H}(\rho[c]_w), t) \in \verdicts, t < \floor{\nicefrac{2}{3}\mathsf{V}} \} \\
        \rho[c] \!\!\!\!\otherwise
    \end{cases}\!\!\!\!\!\!\!
}
\newcommand{\equationtenpointsixteen}{\goodset' \equiv \goodset \cup \{r \mid \tup{r, \floor{\nicefrac{2}{3}\mathsf{V}} + 1} \in \verdicts \}}
\newcommand{\equationtenpointseventeen}{\badset' \equiv \badset \cup \{r \mid \tup{r, 0} \in \verdicts \}}
\newcommand{\equationtenpointeighteen}{\wonkyset' \equiv \wonkyset \cup \{r \mid \tup{r, \floor{\nicefrac{1}{3}\mathsf{V}}} \in \verdicts \}}
\newcommand{\equationtenpointnineteen}{\offenders' \equiv \offenders \cup \{ k \mid (r, k, s) \in \xtculprits \}\cup \{ k \mid (r, v, k, s) \in \xtfaults \}}
\newcommand{\equationtenpointtwenty}{
    \begin{align}
        \mathbf{H}_o \equiv \sq{ k \mid (r, k, s) \in \xtculprits} \frown \sq{ k \mid (r, v, k, s) \in \xtfaults }
    \end{align}
}
